\chapter{Roofing, Cladding and Facade Risk Assessment}
\label{chap:Roofing, Cladding and Facade Risk Assessment}

%---------------------------- SECTION ----------------------------
\section{Why this assessment matters in construction}

\paragraph{Roofing, cladding, and façade operations consistently feature in the HSE's annual fatal injury statistics for the construction industry. Working at the perimeter and on the surface of a building at height, handling large and heavy panels, working in weather conditions that cannot be controlled, and operating adjacent to open roof edges and roof lights create a hazard profile that is among the most unforgiving in the built environment sector. Falls from roofs and through fragile roof surfaces remain among the most common causes of construction fatalities, and the Working at Height Regulations 2005 (WAHR 2005) were substantially shaped by the evidence of these recurring incidents. A roofer or cladding operative working on the external envelope of a building is exposed to the full consequences of a fall --- typically at a height where the fall will be fatal --- with no second chance and no margin for a poorly selected or improperly rigged fall-prevention system.}

\paragraph{The fragility of roof surfaces is a particularly persistent hazard. The UK's existing building stock contains large areas of fibre cement (formerly asbestos cement) sheeting, profiled steel sheeting corroded to below safe load-bearing capacity, patent glazing, polycarbonate rooflights, and glass-fibre reinforced panels, all of which are incapable of supporting the weight of a person and will give way without warning. Falls through fragile roof surfaces have killed experienced roofers who stepped onto a panel believing it to be load-bearing. Advisory fall-through warning notices and colour-coded roof safety information systems (such as those promoted by the National Federation of Roofing Contractors) are important but must be supplemented by physical collective fall prevention measures --- boarding, nets below fragile surfaces, and barriers at roof edge --- that do not rely on the operative's identification and avoidance of fragile materials.}

\paragraph{Cladding and façade installation operations introduce specific hazards beyond the roof edge fall risk. Large composite panels, curtain walling units, stone cladding panels, and glass units are heavy, wind-sensitive, and geometrically awkward to handle. Crane-assisted cladding panel installation requires the concurrent management of lifting operations (see Chapter~\ref{chap:Hot Works and Fire Risk Assessment}), the risk of the panel acting as a wind vane and swinging beyond the exclusion zone, the risk of fastener failure dropping a panel during installation, and the risk to operatives working at the panel's final location who must position and connect the unit while it is suspended on the crane hook. Manual handling of cladding components at height combines the risk of musculoskeletal injury with the risk of dropping components onto persons below or losing balance during the handling operation. The sequential nature of cladding installation, in which successive panels cannot be installed until the preceding panel is fixed, creates programme pressure that can compromise the attention given to working safely on partially-installed systems with gaps in the protective cladding skin.}

\barquote{``The roof edge and the fragile roof surface are the same hazard in different guises: both will result in a fall that is, at working height, fatal. The roofer who trusts a fragile panel with their weight, and the operative who leans over an unprotected parapet, are making the same error of judgement. Neither will survive it.''}

\example{Falls Through Fragile Rooflights}{A roofing contractor is engaged to replace the felt membrane on a flat roof over a warehouse building constructed in the 1970s. The roof has a number of patent glazed rooflights of unknown age and condition. The roofing gang arrives on site with an instruction to replace the membrane but without a specific survey of the rooflight condition. Shortly after starting work on the roof surface, one operative steps onto a rooflight assuming it to be solid; the patent glazing, having lost structural integrity through decades of UV degradation, fails immediately and the operative falls 5.5 metres to the warehouse floor, sustaining fatal head injuries. Investigation reveals that: no pre-work survey of roof light condition had been commissioned; the risk assessment did not specifically address rooflight fragility; no boards had been laid to bridge across the rooflights; and no nets or soft landing systems had been installed below the rooflights. A perimeter edge protection system had been installed around the roof perimeter but offered no protection for the internal rooflight fall hazard. The employer is prosecuted under WAHR 2005 Regulation~6 (failure to prevent falls through fragile surfaces) and receives a custodial sentence suspended for two years and a substantial fine.}

%---------------------------- SECTION ----------------------------
\section{Legal and professional framework}

\paragraph{The primary legal framework for roofing, cladding, and façade risk assessment is the Working at Height Regulations 2005, which impose a hierarchy of duties --- avoid work at height where possible, prevent falls where work at height is necessary, minimise the distance and consequence of a fall where prevention is not reasonably practicable --- and specifically address work near fragile surfaces (Regulation~9), the selection and use of work equipment for work at height (Regulations~5--8), and the inspection of work equipment used at height (Regulation~12). Key frameworks relevant to this chapter include: Health and Safety at Work etc.\ Act 1974; Management of Health and Safety at Work Regulations 1999; Working at Height Regulations 2005; Lifting Operations and Lifting Equipment Regulations 1998 (LOLER); Provision and Use of Work Equipment Regulations 1998 (PUWER); Construction (Design and Management) Regulations 2015; Manual Handling Operations Regulations 1992; Control of Asbestos Regulations 2012 (for legacy roofing materials); and HSE guidance document HSG33 (Health and Safety in Roof Work).}

\begin{itemize}
  \bitem{Working at Height Regulations 2005}{The primary regulatory instrument for all work at height including roofing and cladding. Regulation~6 specifically requires that, where work is to be carried out on or near a fragile surface, suitable and sufficient steps are taken to prevent any person falling through the surface. This duty is not discharged by warning notices alone: physical measures --- boards, nets, or barriers --- must be provided. Regulation~7 requires that where the distance of a fall cannot be prevented, collective measures to arrest the fall (nets, airbags, safety nets) must be used in preference to personal fall arrest equipment, and PPE is the last resort.}
  \bitem{HSG33 --- Health and Safety in Roof Work (HSE, fourth edition 2012)}{The primary HSE guidance document for roofing operations. HSG33 provides detailed guidance on: the use of edge protection systems for flat and pitched roofs; the selection and use of safety nets, soft landing systems, and running-line systems for pitched roof work; the assessment and management of fragile roof surfaces; the use of mobile elevated work platforms (MEWPs) for access; and the management of temporary roof works including weather protection. HSG33 is the definitive operational reference for roofing risk assessments and will be the primary standard applied in any HSE investigation following a roofing incident.}
  \bitem{LOLER 1998}{Applies to all crane and MEWP operations used in cladding and façade installation, including the crane-assisted erection of curtain walling panels, the use of goods hoists and mast climbers, and the use of building maintenance units for façade access. All lifting operations must be planned by a competent person, supervised, and carried out safely.}
  \bitem{Control of Asbestos Regulations 2012 (CAR 2012)}{Apply where roofing works involve asbestos cement sheets (commonly found as corrugated profiled sheeting on agricultural, industrial, and commercial buildings erected before the year 2000). A refurbishment and demolition asbestos survey (see Chapter~\ref{chap:Asbestos Risk Assessment}) must be commissioned before any roofing works that involve the removal, breaking, or disturbance of potential ACMs. Licensed removal may be required for some asbestos cement work depending on fibre type and condition.}
  \bitem{Manual Handling Operations Regulations 1992}{Apply to the manual handling of cladding panels, roofing sheets, and associated fixings and materials at height. A manual handling assessment must address the weight and geometry of panels, the working position required for installation (often at the limit of reach above head height), and the effect of personal fall arrest equipment on freedom of movement and handling capacity.}
  \bitem{CDM 2015}{Requires that roofing and cladding risks are considered in the pre-construction information (particularly the location of fragile surfaces, rooflights, and edge conditions), that the Construction Phase Plan addresses edge protection, fragile surfaces, and the management of weather during roofing works, and that the Health and Safety File records the location of rooflights, fragile surfaces, and fall arrest anchor points for future maintenance work.}
\end{itemize}

%---------------------------- SECTION ----------------------------
\section{Conducting the assessment using the five-step method}

\subsection{Step 1 --- Identify Hazards}
\paragraph{Roofing, cladding, and façade hazards fall into four categories: falls from height (roof edge, roof opening, fragile surfaces, access equipment); struck-by hazards (falling materials, wind-borne panels); manual handling and musculoskeletal hazards; and weather-related hazards (wet surfaces, wind loading on large panels, sun exposure). Falls from height at the roof edge include: falls over an unprotected parapet or eaves; falls through an open roof light or service penetration; falls from a ladder or scaffold used for roof access; and falls from a MEWP or goods hoist basket. Falls through fragile surfaces include: all heritage patent glazing and polycarbonate rooflights regardless of apparent condition; corrugated fibre cement sheets (which lose structural integrity with age and UV exposure and cannot safely support a person even when recently installed if the sheet is cold or has surface condensation); GRP dome rooflights; thin glass roofs; and any profiled sheeting that shows signs of corrosion, crazing, or delamination. Struck-by hazards include: tools or fixings dropped from the roof level onto persons below; lifting slings or crane loads that swing beyond the exclusion zone during cladding installation; and sheet materials that become wind-borne during installation before being fixed. Manual handling hazards include: handling of large heavy panels at the limit of reach; repetitive fixing operations in awkward overhead positions; and carrying materials up or across a pitched roof surface.}

\subsection{Step 2 --- Decide Who or What Is Affected}
\paragraph{Persons at risk from roofing, cladding, and façade hazards include the roofers and cladding operatives working at height; crane drivers and slingers involved in panel lifts; scaffolders maintaining the access scaffold; operatives of other trades who may pass below the working area during materials handling or panel installation; members of the public on adjacent footways who may be struck by falling materials; and future maintenance workers who will access the completed roof --- the CDM 2015 duty to record the location of fall arrest anchors and fragile areas in the Health and Safety File is specifically aimed at protecting these future users.}

\subsection{Step 3 --- Evaluate Likelihood and Consequence}
\paragraph{A fall from a roof edge at typical working heights of 4 to 20 metres carries a consequence of fatal or life-changing injury; the consequence rating is Catastrophic. The likelihood of a fall in the absence of adequate edge protection depends on the duration of the work, the slope of the roof, the weather conditions, and the experience and attentiveness of the operatives. The combination of Catastrophic consequence and material likelihood places roof edge falls and falls through fragile surfaces among the highest overall risk ratings on any construction project. The assessment must distinguish between the controlled risk of a compliant system with adequate collective fall prevention and the uncontrolled risk of work on a fragile or unprotected roof.}

\subsection{Step 4 --- Select and Implement Controls}
\paragraph{Controls must begin with the hierarchy in WAHR 2005: avoid the work at height where possible; prevent falls using collective measures (edge protection, covers for openings, boards over fragile surfaces); and only as a last resort use personal fall arrest equipment. Edge protection must comply with the specification in HSG33 and BS EN 13374 and must be installed and inspected by a competent person before roofing work commences. All rooflights and openings must be covered or guarded before any operative works near them. Where fragile surfaces cannot be completely covered, safety nets or soft landing systems must be installed below them before work commences. Personal fall arrest systems using running-line anchors may be used for steep pitch roof work where the above measures are impracticable, but only where anchors are certified, the lines and harnesses are inspected, and the operatives are trained in the use of their equipment.}

\subsection{Step 5 --- Record and Review}
\paragraph{The roofing risk assessment must be reviewed at the start of each day on which roofing work is planned, taking account of weather conditions. Work must not commence if wind speeds exceed the limits specified in the risk assessment for the specific activities planned --- HSG33 provides guidance on wind speed limits for various roofing and cladding activities. The inspection records for edge protection, safety net installation, and personal fall protection equipment must be maintained. Any incident, near miss, or change in weather or site conditions that affects the risk profile must trigger an immediate review.}

\example{Five-Step Application}{A roofing contractor is replacing the membrane on a 500 m\textsuperscript{2} flat roof on a five-storey commercial building. The roof has twelve patent glazed rooflights of 1960s vintage. \textbf{Step~1 (Hazards):} Falls from unprotected roof edge (4.5 m to hardstanding); falls through patent glazed rooflights (fragile); falls through existing bituminous felt (potential delamination); material dropped onto public footway adjacent to building. \textbf{Step~2 (Who Affected):} Roofers, scaffolders, public on adjacent pavement. \textbf{Step~3 (Evaluation):} Falls rated Catastrophic/High (inherent, no edge protection); fragile surface rated Catastrophic/High. \textbf{Step~4 (Controls):} Tube-and-fitting scaffold with double-guardrail and toe-board erected by competent scaffolder before roofers commence; all 12 rooflights boarded over with 18~mm plywood secured in place before any roofer walks on roof; protection fans over public footway; materials hoist for materials delivery to roof level; materials secured against wind before end of each working day. \textbf{Step~5 (Record):} Weekly scaffold inspection records; daily check of rooflight boarding before work commences; pre-work weather check (work suspended above Beaufort 5).}

%---------------------------- SECTION ----------------------------
\section{Applying the hierarchy of controls}

\paragraph{The hierarchy of controls for roofing and cladding follows WAHR 2005 precisely: avoid work at height; use collective fall prevention measures; use work restraint; use collective fall arrest; use personal fall arrest (harness and lanyard) only as the last resort. At each level, the specific equipment selected must be appropriate for the roof type, the work activity, and the skill level of the operatives.}

\begin{itemize}
  \bitem{Elimination}{Consider whether the proposed roofing or cladding work can be designed to be carried out without requiring personnel to be exposed to the fall risk. Prefabricated roof modules assembled at ground level and craned into position eliminate the need for operatives to work on the roof surface during assembly. Cladding panels designed for crane-only installation with fixing connections that can be made without operatives leaning out of a MEWP reduce the exposure at height. Designers specifying roof and cladding systems have a CDM 2015 duty to consider and, where practicable, design out working at height risks.}
  \bitem{Substitution}{Where collective fall prevention cannot be provided from scaffold, substitute tower scaffold or a MEWP for working from the roof surface. A MEWP with a work restraint system keeps the operative within the basket and eliminates the roof-edge fall risk for overhanging or perimeter fixing operations, provided the MEWP can be positioned appropriately.}
  \bitem{Engineering controls}{Edge protection to BS EN 13374 (guardrail, mid-rail, and toe-board) is the primary engineering control for flat and shallow-pitch roofs. Safety nets installed below the working area provide collective fall arrest without constraining the operatives' movement on the roof. Boarding and covers for fragile surfaces prevent falls through the surface. Exclusion zones below the working area prevent struck-by injuries from falling materials. Hoards, fans, or nets over public footways protect members of the public.}
  \bitem{Administrative controls}{Pre-work weather check and defined wind speed limits for cessation of work. Daily briefing on the specific activities and controls for the day's work. Permit-to-proceed system for roofing operations on days when conditions are marginal. Prohibition on working over fragile surfaces without boards in place. Training requirements for all operatives in the use of personal fall arrest systems. CDM record of fragile surfaces and fall arrest anchors for future maintenance.}
  \bitem{PPE}{Personal fall arrest harness and lanyard, or work restraint system, selected and inspected in accordance with PUWER and WAHR 2005. Hard hat with chin strap (to prevent it being knocked off at height). Non-slip footwear. High-visibility vest. Respiratory protection where asbestos cement is being disturbed (see Chapter~\ref{chap:Asbestos Risk Assessment}).}
\end{itemize}

%---------------------------- SECTION ----------------------------
\section{Cadence of assessment and review triggers}

\paragraph{Roofing and cladding risk assessments must be reviewed at the start of each day of operations and following specific trigger events that change the risk profile.}

\begin{itemize}
  \bitem{Before commencement of each day's work}{Wind speed and weather conditions must be assessed before work commences each day. The risk assessment must specify the wind speed at which different activities must cease, and the pre-work check must verify that edge protection and fragile surface covers are in place and undamaged.}
  \bitem{Following adverse weather overnight}{High winds, frost, or heavy rain may damage edge protection, displace fragile surface covers, and make the roof surface or access equipment slippery. A pre-work inspection after adverse weather is mandatory before any operative goes onto the roof.}
  \bitem{When a new phase of works exposes a new roof area}{As roofing progresses across a large roof area, previously protected areas may become unprotected as operations move, or new fragile surfaces may be exposed. The risk assessment must be reviewed at each stage of progression across the roof.}
  \bitem{Following any incident or near miss}{Any fall arrest activation, near miss, or incident on the roof must trigger an immediate formal review of the controls in place and the specific circumstances that led to the event.}
  \bitem{When the work scope changes}{If additional areas of roof are added to the scope, if the roof construction is found to differ from the pre-work survey (additional rooflights, a previously unknown fragile section), or if weather delays require work to be continued into conditions outside the original risk assessment parameters, the assessment must be reviewed before work continues.}
\end{itemize}

%---------------------------- SECTION ----------------------------
\section{Common failure modes}

\begin{itemize}
  \bitem{Edge protection not installed before roofers start work}{The most common fatal scenario in roofing is the fall from an unprotected edge at the commencement of a project, before edge protection is installed. Roofers attend site, the scaffold is not yet erected or has been stripped, and work begins without edge protection. The risk assessment's requirement for edge protection to be in place before roofing commences must be enforced as an absolute requirement, not subject to informal override.}
  \bitem{Fragile rooflights not identified or covered before roof access}{Pre-work surveys that miss rooflights --- due to inaccurate building plans, obscured surfaces, or the operative conducting the survey not being competent to identify fragile materials --- result in roofers working on a surface that contains fatal fall hazards they do not know are there. A competent pre-work survey of the roof condition, specifically to identify all fragile surfaces, is a prerequisite for all roofing works on existing buildings.}
  \bitem{Personal fall arrest equipment used incorrectly}{Personal fall arrest harnesses used with lanyards that are too long (allowing a fall distance greater than the available clearance below the operative), with anchor points not rated for dynamic arrest load, or without adequate training in donning and doffing, provide false security while delivering no effective fall arrest capability. The assessment must specify the maximum lanyard length for the available working height above the fall risk.}
  \bitem{Work continued in unsafe wind conditions}{Programme pressure to complete roofing works before weather windows close creates a tendency to continue work in wind conditions beyond the safe limit. Sheet materials, boards, and partially-installed panels that are not secured in high wind become projectiles or sails. Work in dangerous wind conditions has killed site workers adjacent to building facades.}
  \bitem{No exclusion zone below working area}{Roofing and cladding operations frequently fail to establish and enforce an exclusion zone at the base of the building directly below the working area. Dropped tools, loose fixings, sheet offcuts, and panels dropped during installation all represent fatal struck-by hazards for persons below. Exclusion zones must be established and enforced for the full duration of roofing and cladding operations.}
\end{itemize}

\example{Failure Pattern}{A roofing contractor is replacing profiled metal sheeting on the roof of an industrial building. The work requires access to the full width of the roof, 12 metres above the internal floor level. Safety nets have been specified in the risk assessment but have not been installed due to cost and programme discussions that were never resolved. The contractor proceeds with edge protection (guardrails) at the eaves only. In the course of the work, an operative steps onto a corrugated fibre cement sheet covering a rooflight position that was not identified on the building owner's plans. The sheet fails and the operative falls 12 metres to the concrete floor below, sustaining fatal injuries. Post-incident survey identifies that the building has twelve further unrecorded rooflights of the same material. Investigation confirms that: no pre-work fragile surface survey had been carried out beyond a visual check from the ground; the safety nets specified in the risk assessment had not been installed; the risk assessment had been reviewed and approved without verifying that the nets were in place; and the principal contractor had signed off a stage completion payment without checking compliance with the fall prevention requirements.}

%---------------------------- CHAPTER SUMMARY ----------------------------
\newpage
\section{Chapter Summary}

\begin{table}[H]
\centering
\begin{tabularx}{\textwidth}{>{\raggedright\arraybackslash}p{3.2cm}X}
\toprule
\textbf{Assessment Element} & \textbf{Operational Requirement} \\
\midrule
Legal/Professional Basis & Working at Height Regulations 2005; CDM 2015; LOLER 1998; PUWER 1998; CAR 2012 (for ACM roofing); Manual Handling Regulations 1992; HSG33; Health and Safety at Work etc.\ Act 1974. \\
Method & Apply the full five-step process and document inherent/residual risk. Conduct a pre-work fragile surface survey for all works on existing buildings. Apply the WAHR 2005 hierarchy: prevent falls first with collective measures. \\
Controls & Edge protection to BS EN 13374 installed before roofing commences; all fragile surfaces boarded or netted; exclusion zone below working area; no work above Beaufort 5 for panel installation; daily pre-work inspection of edge protection and fragile surface covers. \\
Cadence & Daily pre-work check; after adverse weather; when new roof area exposed; following any incident or near miss; when scope changes. \\
Assurance & Weekly scaffold/edge protection inspection records; pre-work weather log; fragile surface survey report; CDM Health and Safety File records of fragile surfaces and anchor points. \\
\bottomrule
\end{tabularx}
\end{table}

\begin{itemize}
  \bitem{Key takeaway 1}{Edge protection to the standard required by HSG33 and BS EN 13374 must be in place before any roofer or cladding operative commences work on the roof. This is an absolute requirement, not subject to informal override for programme or access reasons.}
  \bitem{Key takeaway 2}{A pre-work survey of all fragile surfaces --- conducted by a competent person with physical access to the roof --- is a prerequisite for all works on existing building roofs. All fragile surfaces must be boarded or netted before any operative walks on the roof.}
  \bitem{Key takeaway 3}{The WAHR 2005 hierarchy must be applied: prevent falls with collective measures (edge protection, nets, boarding) before considering personal fall arrest equipment. Harnesses and lanyards are the last resort, not the primary control.}
  \bitem{Key takeaway 4}{An exclusion zone at the base of the building, below the full width of the working area, must be established and enforced for the duration of all roofing and cladding operations. Dropped materials from roof level are a fatal hazard for persons at ground level.}
  \bitem{Key takeaway 5}{Roofing and cladding risk assessments must specify defined wind speed limits for the cessation of sheet and panel installation operations, and work must be suspended when those limits are exceeded regardless of programme pressure.}
\end{itemize}

\barquote{In Chapter~\ref{chap:Occupational Health and Dust Risk Assessment}, we examine the assessment and control of occupational health hazards arising from dust, fume, and other airborne contaminants generated by construction activities --- hazards that kill more construction workers each year through occupational disease than all acute traumatic incidents combined.}
